% !PRINTED 240
\weberpage{240}
%<<CONT
durch Composition von $D_2$ mit
\[
j_1 = \begin{pmatrix} 0, & -\sqrt{i} \\ i\sqrt{i}, & 0 \end{pmatrix}
\]
entsteht, also
\[
D_2' = D_2 + D_2 j_1
\]
ergiebt. Nun ist allgemein:
\[
j_1^{-1} \begin{pmatrix} \alpha, & \beta \\ \gamma, & \delta \end{pmatrix} j_1 = \begin{pmatrix} \delta, & \gamma i \\ -\beta i, & \alpha \end{pmatrix},
\]
und daraus schliesst man nach \S.\,56, (4), dass $j_1^{-1} T j_1 = T$ ist, dass also aus der Tetra\"edergruppe zwei erweiterte Gruppen $T' = T+Tj$ und $T'' = T+Tj_1$ abgeleitet werden können; aber auch nur zwei, weil nach \S.\,56 die Tetra\"edergruppe durch eine darin enthaltene Vierergruppe völlig bestimmt ist. Denn wenn wir statt der Erweiterung $D_2'$ die Erweiterung der Vierergruppe zu einer der Gruppen (11) wählen, so erhalten wir eine erweiterte Tetra\"edergruppe $T'''$, die durch Transformation mit $\varphi_1$ oder $\varphi_2$ in $T''$ übergeht.

Da die Substitution
\[
\begin{pmatrix} 0, & e^{\frac{\pi i}{8}} \\ e^{-\frac{\pi i}{8}}, & 0 \end{pmatrix}
\]
mit der Octa\"edergruppe \S.\,57, (8) nicht vertauschbar ist, so lässt sich aus der Octa\"edergruppe ausser durch Inversion keine weitere Gruppe zweiter Art ableiten, und dasselbe gilt von der Ikosa\"edergruppe. Man muss, um diese Schlussweise anzuwenden, die Ikosa\"edergruppe so transformiren, dass eine der darin enthaltenen Vierergruppen die einfachste Gestalt annimmt, was etwas weitläufig, aber durchaus nicht schwierig ist, und hier nicht weiter ausgeführt werden soll.

Wir erhalten also noch folgende Gruppen zweiter Art:
\[
\begin{aligned}
&\text{IV.}\quad 1)\ T' = T+Tj, \qquad 2)\ T'' = T+Tj_1.\\
&\text{V.}\qquad\qquad O' = O+Oj.\\
&\text{IV.}\qquad\qquad J' = J+Jj.
\end{aligned}
\]
% [sic] the third line is printed "IV." again (as the first); expected "VI." given the sequence — reproduced as printed.
